\documentclass[11pt]{article}
% Page layout and typography
\usepackage[margin=1in]{geometry}
\usepackage{microtype}
\usepackage{parskip}

% Mathematics
\usepackage{amsmath,amssymb,amsthm,mathtools}
\usepackage{bm}

% Figures and links
\usepackage{graphicx}
\graphicspath{{figures/}}
\usepackage{xcolor}
\usepackage[hidelinks]{hyperref}

% Lists
\usepackage{enumitem}
\setlist{nosep}

% Theorem environments
\theoremstyle{definition}
\newtheorem{definition}{Definition}[section]
\newtheorem{example}[definition]{Example}
\newtheorem{remark}[definition]{Remark}

\theoremstyle{plain}
\newtheorem{theorem}[definition]{Theorem}
\newtheorem{lemma}[definition]{Lemma}
\newtheorem{proposition}[definition]{Proposition}
\newtheorem{corollary}[definition]{Corollary}

% Common number systems
\newcommand{\R}{\mathbb{R}}
\newcommand{\C}{\mathbb{C}}
\newcommand{\Q}{\mathbb{Q}}
\newcommand{\Z}{\mathbb{Z}}
\newcommand{\N}{\mathbb{N}}
\newcommand{\D}{\mathbb{D}}
\newcommand{\T}{\mathbb{T}}
\newcommand{\RS}{\mathbb{S}^2}
\newcommand{\per}{\text{per}}

% Common operators
\DeclareMathOperator{\im}{im}
\DeclareMathOperator{\rank}{rank}
\DeclareMathOperator{\Span}{span}
\DeclareMathOperator{\id}{id}


\title{STAT5351 Lecture Notes}
\author{Caelen Gelineau}
\date{\today}

\begin{document}

\maketitle
\tableofcontents
\newpage

\section{Lecture 1 - 24/08/26}
\subsection{Counting principle}
Suppose you can solve a problem in 2 steps, with $n_1$ ways for step 1 and $n_2$ ways for step 2. \\
How many ways can you solve the problem?
\[
    \begin{matrix}
        (x_1, y_1) & (x_1, y_2) & \dots & (x_1, y_{n_2})\\
        \vdots\\
        (x_{n_1}, y_1) & (x_{n_1}, y_2) & \dots & (x_{n_1}, y_{n_2})
    \end{matrix}
\]
Thus, there are $n_1 \cdot n_2$ ways.
\begin{theorem}
    Given $r$ steps to solve a problem and $n_i$ ways to solve step $i$, then the problem can be
    solve in $n_1 \cdot n_2 \cdot \dots \cdot n_r$ ways.
\end{theorem}
\subsection{Permutations}
\begin{definition}
    A permutation on $n$ different elements is a $\textbf{DISTINCT}$ arrangement of the $n$ elements, where distinct arrangement
    means no repetition of elements in the arrangement.
\end{definition}

\begin{theorem}
    The number of permutations of $n$ unique elements is $n \cdot (n-1) \cdot (n-2) \dots \cdot 2 \cdot 1 = n!$. 
\end{theorem}

\begin{theorem}
    The number of permutations of $n$ elements taken $r$ at a time is \\
    $\text{nPr} = n(n-1)(n-2)\cdots (n-r+1) = \frac{n!}{(n-r)!}$
\end{theorem}

\begin{theorem}
    Suppose $n$ elements have $k$ different types: $n_1$ of type 1 to $n_k$ of type k. Then $n = n_1 + \cdots + n_k$ and the number
    of permutations of $n$ elements is $\frac{n!}{n_1!n_2!\cdots n_k!}$.
\end{theorem}

\subsection{Combinations}
\begin{definition}
    A combination of $n$ different elements taken $r$ at a time is a selection of $r$ elements from $n$ elements without regard to the order of selection.
\end{definition}
\begin{theorem}
    The number of permutations of $n$ elements taken $r$ at a time ($n$ choose $r$) is 
    \[
        \text{nPr} = \frac{n!}{(n-r)!r!}
    \]
    
\end{theorem}

\section{Lecture 2 - 26/08/26}
\subsection{Partitions}
\begin{definition}
    A partition of $n$ elements into $k$ subsets is a way to separate $n$ elements into $k$ groups without regard for order.
\end{definition}

\begin{theorem}
    The number of ways in which $n$ elements can be partitioned into $k$ subsets such that $n_1$ elements in subset 1, 
    $n_2$ in subset 2, etc is 
    \[
        \frac{n!}{n_1!n_2! \cdots n_k!}
    \]
\end{theorem}

\section{Lecture 3 - 31/08/26}
\subsection{Discrete probability function}
Discrete probability function $P$ has the following properties:
\begin{enumerate}
    \item $P(S) = 1$.
    \item $P(\emptyset) = 0$.
    \item $P(A^c) = 1 - P(A)$.
    \item $P(A \cup B) = P(A) + P(B) - P(A \cap B)$
    \item If $A \cap B = \emptyset$, $P(A \cup B) = P(A) + P(B)$.
    \item $A, B \subset S$, $A \subseteq B \implies P(A) \le P(B)$.
\end{enumerate}

\section{Lecture 4 - 02/09/26}
\subsection{Conditional Probability}

If $B$ has occurred, what is the probability of $A$?
\begin{theorem}
    $P(A|B) = \frac{P(A \cap B)}{P(B)}$
\end{theorem}
\begin{theorem}
    $P(A \cap B) = P(A|B) \cdot P(B)$
\end{theorem}
\begin{theorem}
    $P(A \cap B \cap C) = P(A)P(B|A)P(C|A \cap B)$
\end{theorem}

\section{Lecture 5 - 09/09/26}
\subsection{Bayes' Theorem}
\begin{definition}
    $A_1, \dots, A_k \subset B$ form a partition of B if $A_i$ are all mutually exclusive.
\end{definition}
\begin{remark}
    For a partiton of $B$,
    \[
        P(B) = P(A_1 \cap B) + \cdots + P(A_k\cap B) = P(B|A_1)P(A_1) \cdots P(B|A_k)P(A_k)
    \]
\end{remark}
\begin{theorem}
    $P(B) = \sum_{n=1}^{k}P(B|A_i)P(A_i)$.
\end{theorem}
\begin{theorem}[Bayes' Theorem]
    \[
        P(A_i | B) = \frac{P(A_i)P(B|A_i)}{\sum_{j} P(B|A_j)P(A_j)}
    \]
\end{theorem}
Bayesian inference is forming an event $B$ to conclude the cause, i.e. which $A_i$ is the cause of $B$.
\subsection{Independent Events}
\begin{definition}
    Two events are independent if 
    \[
        P(A \cap B) = P(A)P(B)
    \]
    if they are not independent they are dependent.
\end{definition}
\begin{theorem}
    If $A$ and $B$ are independent, then $A$ and $B^c$ are independent.
\end{theorem}

\section{Lecture 6 - 14/09/26}
\subsection{Random Variables}
\begin{definition}
    Let $S$ be a sample space with a probability function $P$. A random variable is a function $X: S \rightarrow \R$.
    The notation is $P(X=b) :=$ probability of event $\{X = b\}$.
\end{definition}

\section{Lecture 7 - 16/09/26}
\subsection{Probability mass function}
\begin{theorem}
Let $X$ be a discrete random variable on $S$ such that $b_1, \dots, b_n$ are the values in the image of $X$. Then $\{X=b_1\}, \dots, \{X =b_n\}$ form a partition of $S$.
\end{theorem}
\begin{definition}
    The probability mass function of $X$ is 
    \[
        f(y) := P(X = y), \quad \forall y \in \text{image of } X
    \]
\end{definition}
\begin{theorem}
    If $f(y)$ is a pmf of $X$, then
    \begin{enumerate}
        \item $f(y) \ge 0$.
        \item $\sum_{y}f(y) = 1 \iff \sum f(b_i) = 1$.
    \end{enumerate}
\end{theorem}
\subsection{Cumulative Distribution}
\begin{definition}
    $F(y):= P(X \ge y), \quad \forall y \in -\infty < y < \infty$ is the cumulative distribution of $X$. It measures the accumulation of probability as $y$ increases.
\end{definition}

\section{Lecture 8 - 21/09/26}
\begin{example}
$F(y)$ is the cumulative distribution of X.
\[
    F(y) = 
    \begin{cases}
        0 \quad y < 1\\
        \frac{1}{3} \quad 1 \le y < 4\\
        \frac{1}{2} \quad 1 \le 4 < 6 \\
        \frac{5}{6} \quad 6 \le 4 < 10 \\
        1 \quad y \ge 10 
    \end{cases}
\]
$P(X = 4) =$ ?, $P(2 < X \le 6) = $ ?\\
First, find the probability mass function (p.m.f). X takes values $b_1 = 1, b_2 = 4, b_3 = 6, b_4 = 10$.\\
Recall that $F(b_{i + 1}) = F(b_i) + f(b_i+1)$ where f is the p.m.f.\\
\[
    f(b_1) = f(1) = F(1) - 0 = \frac{1}{3}
\] 
\[
    f(b_2) = f(4) = F(4) - F(1) = \frac{1}{2}- \frac{1}{3} = \frac{1}{6}
\]
\[
    f(b_3) = \frac{1}{3}
\]
\[
    f(b_4) = \frac{1}{6}
\]
Thus,
\[
    P(X=4) = f(4) = \frac{1}{6}
\]
and since 4 is inside the interval $(2, 6]$, then 
\[
    (X=4) \cup (X=6)
\]
and since the events are mutually exclusive,
\[
    P(2 < X \le 6) = P(X = 4) + P(X=6)
\]
\[
    = f(4) + f(6) = \frac{1}{2}
\]
\end{example}
\begin{example}
$X$ is a discrete random variable with cumulative distribution
\[
    F(x) =
    \begin{cases}
        0 \quad y < -1\\
        \frac{1}{9} \quad -1 \le y < 1\\
        \frac{5}{9} \quad 1 \le y < 3 \\
        \frac{8}{9} \quad 3 \le y < 5\\
        1 \quad y \ge 5
    \end{cases}
\]

Find the pmf of $X$ and $P(0 \le X < 5)$.\\
$X$ takes values $-1, 1, 3, 5$. Find the pmf:
\[
    f(-1) = \frac{1}{9}
\]
\[
    f(1) = \frac{5}{9} - \frac{1}{9} = \frac{4}{9}
\]
\[
    f(3) = \frac{1}{3}
\]
\[
    f(5) = \frac{1}{9}
\]
\[
    P(0 \le X < 5) = P(X=1) + P(X=3) = f(1) + f(3) = \frac{7}{9}
\]
\end{example}
\subsection{Continuous random variables}
\begin{definition}
    Let $S$ be a continous sample space. $X: S \rightarrow \R$ is a continous random variable if there 
    exists a function $f(t) \ge 0$ defined on $\R$ such that 
    \[
        P(a \le X \le b) = \int_{a}^{b}f(t) dt
    \]
    $f(t)$ is called the probability density function of $X$.
\end{definition}
\begin{remark}
    $f(t) \ge 0$ is defined for ALL $t \in \R$.
\end{remark}
\begin{theorem}
    Let $X$ be a continuous random variable, and $f$ the pdf of $X$. Then
    \begin{itemize}
        \item $P(X = a) = 0$
        \item $P(a \le X \le b) = P(a \le X < b) = \dots$
        \item $\int_{-\infty}^{\infty} f(t)\, dt$ = 1
    \end{itemize}
\end{theorem}

\begin{definition}
    Let $X$ be a continous random variable and $f(t)$ the pdf. THen
    \[
        F(y) = P(X \le y) = \int_{-\infty}^{y}f(t) \, dt
    \]
    is called the cumulative distribution function of X.
\end{definition}

\begin{theorem}
   Let $f$ be the pdf and $F$ the cdf of X.
   \begin{itemize}
    \item $P(a \le X \le b) = F(b) = F(a)$
    \item if $f$ is a continuous function on $\R$ then $f(y) = \frac{dF(y)}{dy}$.
   \end{itemize}
    
\end{theorem}

\subsection{\texorpdfstring{Comparison of discrete and continous random variables $X$}{Comparison of discrete and continous X}}
\subsubsection{Discrete}
\begin{itemize}
    \item $f$ is pmf
    \item $f$ defined on discrete values
    \item $F$ is a cumulative distribution
    \item $F(y)$ = $\sum_{b_i \le y}^{\infty} f(b_i)$
    \item $P(a < X \le B) = F(b) - F(a)$, the endpoints matter
\end{itemize}
\subsubsection{Continuous}
\begin{itemize}
    \item $f$ is a pdf
    \item $f$ is defined on $\R$
    \item $F$ is a cdf
    \item $F(y) = \int_{-\infty}^{y} f(t)\, dt$
    \item $P(a \le X \le b) = \int_{a}^{b} f(t) \, dt$, the endpoints don't matter.
\end{itemize}

\section{Lecture 9 - 23/09/26}
\subsection{Joint/marginal distribution of discrete random variables}
For discrete random variables $X, Y$, can we find $P(X=Y)$ by using the pmf $f,g$ respectively? 
No, we need more information about their relation.
\begin{definition}
    $X, Y$ disc. rand. vars. 
    \[
        f(x, y) = P(X=x, Y=y)
    \]
    for the pair $(x,y)$ in the range of $X, Y$ is called the joint probability mass function of $X,Y$.
\end{definition}
\begin{theorem}
    if $f(x,y)$ is joint pmf,
    \begin{enumerate}
        \item $f(x,y) > 0$
        \item $\sum_{x} \sum_{y} f(x,y) = 1$
    \end{enumerate}
\end{theorem}

\begin{definition}
    Joint cumulative distribution of discrete random variables $X,Y$ is defined as
    \[
        F(x,y) = P(X \le x, Y \le y) = \sum_{s\le x} \sum_{t \le y}f(s,t)
    \]
\end{definition}
\begin{remark}
    Joint pmf is defined on discrete values but the joint cumulative distribution is defined on $\R^2$.
\end{remark}
Suppose we are given $f(x,y)$, the joint pmf of $X$ and $Y$. How do we find the pmfs of $X$ and $Y$? Idk.
\begin{definition}
    $f(x,y)$ the joint pmf of $X,Y$. For each $x$ in range of $X$, 
    \[
        g(x) = \sum_{y} f(x,y)
    \]
    is called the marginal pmf of $X$. The marginal pmf of $X$ is just the pmf of $X$.
    
\end{definition}
\begin{definition}
    $f(x,y)$ is a joint pdf of $X, Y$. if it satisfies
    \begin{enumerate}
        \item $f(x,y) \ge 0$ it is defined for all $(x,y) \in \R^2$.
        \item $P(X,Y) \in A = \int \int_A f(x,y)dxdy$
    \end{enumerate}
    
\end{definition}

\begin{theorem}
    $f(x,y)$ is a joint pdf. Then
    \[
        \int_{-\infty}^{\infty} \int_{-\infty}^{\infty} f(x,y)\, dxdy = 1
    \]
    
\end{theorem}

\section{Lecture 10 - 28/09/26}
Exam 1 Review.

\section{Lecture 11 - 05/10/26}
\subsection{Marginal pdfs and pmfs}
\begin{definition}
    Let $X,Y$ be continuous random variables. The joint c.d.f. 
    \[
        F(x,y) = P(X \le x, Y \le y) = \int_{-\infty}^{y}\int_{-\infty}^{x} f(s,t)\, ds \, dt
    \]
    where $f(x)$ is the joint p.d.f.
    
\end{definition}
Question: Given the joint pmf of $X, Y$, can we find the pdf of $X$ and $Y$?
\begin{definition}
    $X, Y$ continuous random variables. Let $f(x,y)$ be the joint pdf. The marginal pdf of $X$ is
    \[
        g(x) = \int_{-\infty}^{\infty} f(x,y)\, dy, \quad -\infty < x < \infty
    \]
\end{definition}
\begin{remark}
    The marginal pdfs are just the pdfs of $X$ and $Y$ respectively.
\end{remark}

\begin{definition}
    Let $f(x,y)$ be the joint pmf of discrete random variables $X,Y$. Let $h(y)$ be the marginal pmf of $Y$.
    Then 
    \[
        f(x|y) = P(X=x,Y=y) = \frac{f(x,y)}{h(y)}
    \]
    for all $x \in \text{range}(x)$. $f(x|y)$ is called the conditional pmf of $X$ given $Y=y$. 
\end{definition}
\begin{remark}
    We can regard the conditional probability as a function of two variables.
\end{remark}

\section{Lecture 12 - 07/10/26}
\begin{definition}
    The joint pdf is $f(x,y)$ of continuous random variables $X,Y$. Let $h(y)$ be the marginal pdf of $Y$.
    \[
        f(x|y) = \frac{f(x,y)}{h(y)}
    \]
    is the conditional pdf of $X$ given $Y=y$. We can regard the conditional pdfs as functions of two variables.
\end{definition}
\begin{definition}
    $f(x)$ is the joint pdf of continuous random variables. $g(x)$ is the marginal pdf of $X$. $h(y)$ marginal pdf for $Y$. Then $X,Y$ are independent if 
    \[
        f(x,y) = g(x)h(y), \quad \forall x,y \text{ s.t. }  -\infty < x < \infty, \; -\infty< y < \infty
    \]
    
\end{definition}
\begin{remark}
    If we are given the joint pmf/pdf, then we can find the marginal ones. The converse is not necessarily true. 
    However, if we are given the pdf/pmfs of $X,Y$, then in general we can't find the joint pdf/pmf. When $X,Y$ are independent this is a special case where we can.
\end{remark}

\subsection{Expectation}
After doing some random event many many times, what is the expected value? The "average" of the random variable with respect to the probability distribution.
\begin{definition}
    $X$ is a discrete random variable. $f(x)$ is the pmf. The expectation of $X$ is 
    \[
        E(X) = \sum_{x}x\cdot f(x) \in \R
    \]
    where you sum over all possible values of $X$. For continuous random variable,
    \[
        E(X) = \int_{-\infty}^{\infty} xf(x)\, dx \in \R
    \]
     
\end{definition}

\section{Lecture 13 - 12/10/26}
\subsection{Expectation value for composed functions}
\begin{definition}
    Let $X$ be a random variable on $S$. Define a function
    \[
        g(t): \R \rightarrow \R 
    \]

    Then $g(X)$ is the composition of $g$ and $X$. It is a new random variable on $S$.
\end{definition}

\begin{theorem}
    Let $X$ be a discrete random variable and $f(x)$ is the pmf. Then
    \[
        E(g(X)) = \sum_{x} g(x) \cdot f(x) 
    \]
\end{theorem}

\begin{theorem}
    Let $X$ be a continuous random variable and $f(x)$ is the pdf of $X$. Then
    \[
        E(g(X)) = \int_{-\infty}^{\infty} g(x)f(x)\, dx
    \]
    
\end{theorem}

\begin{theorem}
    Let $a, b \in \R$ be constants. Then
    \[
        E(aX+b) = aE(x) + b
    \]
    
\end{theorem}
\begin{theorem}
    let $c_1, c_2, \dots, c_n$ be constants and $g_i:\R \rightarrow \R$. Then
    \[
        E(\sum_{n=1}^{n}c_ig_i(X)) = \sum_{i}^{n} c_i E(g_i(X))
    \]
    
\end{theorem}

\begin{theorem}
    Let $X, Y$ be discrete random variables and $f(x,y)$ is a joint pmf. then
    \[
        E(g(X,Y)) = \sum_{x} \sum_{y} g(x,y)f(x,y)  
    \]
    
\end{theorem}
\subsection{Moments, mean, standard deviation, variance}
\begin{definition}
    the $r$th moment about the origin of random variable $X$ is $\mu'_r =E(X^r)$.
\end{definition}
\begin{remark}
    For $r = 0$, $\mu'_0 = E(1) = 1$.\\
    For $r = 1$, $\mu'_1 = E(X)$ is called the mean of $X$, denoted by $\mu$.\\
\end{remark}
\begin{definition}
    the $r$th moment about the mean denoted by $\mu_r$ is 
    \[
        \mu_r = E((X-\mu)^r)
    \] 
    where $\mu$ is the mean of $X$.
\end{definition}
\begin{definition}
    Let $X$ be a continuous random variable and $f(x)$ the pdf. Then the mean is given by
    \[
        \mu = \int_{-\infty}^{\infty} xf(x)\, dx
    \]
    and 
    \[
        \mu_r = \int_{-\infty}^{\infty} (x-\mu)^rf(x)\, dx
    \]
\end{definition}
\begin{theorem}
    $\sigma^2 = \mu'_2 - \mu^2 = E(X^2)-E(X)^2$
\end{theorem}
\begin{remark}
    $\sigma^2 = E((x-\mu)^2) = \sum_{x}(x-\mu)^2 f(x)$ for the discrete case and $\int_{-\infty}^{\infty} (x-\mu)^2f(x)\, dx$ for the continuous case.
\end{remark}


\section{Lecture 14 - 14/10/26}
\begin{theorem}
    Let $X$ have variance $\sigma^2$. Consider $aX + b$ where $a, b$ are constant. Then it has 
    \[
        \text{Var}(aX+b) = a^2 \sigma^2
    \]
\end{theorem}
\subsection{Intepretations of standard deviation}
\begin{remark}   
     The standard deviation describes how far the values of $X$ are spread from the mean. Consider the case when $X$ is a discrete random variable.
        \[
            \sigma^2 = E[(X-\mu)^2] = \sum (x-\mu)^2 f(x)
        \]
    where $f(x)$ is the pmf. The pmf is a "weighting", so then the $\sigma^2$ measures the "weighted average" squared distance to the mean. 
    Suppose that $X$ is a continuous random variable, then 
        \[
            P(|x-\mu|\le t) = P(-t \le x-\mu \le t) = P(-t+\mu \le X \le \mu+t) = \int_{\mu-t}^{\mu+t} f(x)\, dx
        \] 
        where $f(x)$ is the pdf. In practice, we may not have all the information about $f$ to calculate the integral above. 
        Since $\sigma^2$ describes the spread of the distribution we can estimate the probability using it.
        This motivates the following theorem:
\end{remark}
\begin{theorem}[Chebyshev's Inequality]
    \[
        P(|x-\mu| \le t) \ge 1 - \frac{\sigma^2}{t^2}
    \]
\end{theorem}
\subsection{Moment generating function}
\begin{definition}
    A moment generating function of a random variable $X$ is 
    \[
        M_X(t) = E(e^{tX})
    \]
    For a discrete random variable,
    \[
        M_X(t) = \sum_{x} e^{tx} f(x)
    \]
    For continuous,
    \[
        M_X(t) = \int_{-\infty}^{\infty} e^{tx}f(x)\, dx
    \]
\end{definition}
\begin{remark} The following properties of $M_X(t)$ hold
    \begin{enumerate}
        \item $M_X(t)$ is a function of $t$.
        \item $M_X(t)$ may not be defined for all $t$.
        \item If $M_X(t)$ is defined, $M_X(t) \ge 0$.
    \end{enumerate}
\end{remark}
Consider the series expansion of $e^{tx}$.
\[
    e^{tx} = 1+tx+\frac{t^2x^2}{2}+\frac{t^3x^3}{3!} + \cdots
\]
Let $X$ be a continuous random variable. then
\begin{align*}
    M_X(t) &= \int_{-\infty}^{\infty} (1+tx+\frac{t^2x^2}{2} + \cdots)f(x)\, dx\\
    &= \int_{-\infty}^{\infty} f(x)\, dx + t\int_{-\infty}^{\infty} xf(x)\, dx + \frac{t^2}{2!}\int_{-\infty}^{\infty}x^2 f(x)\, dx + \cdots \\
    &= \mu_0' + \mu_1't+\mu_2' \frac{t^2}{2!}+\cdots + \mu_r' \frac{t^r}{r!}+\cdots\\
    &= \sum_{r=0}^{\infty} \mu_r' \frac{t^r}{r!}   
\end{align*}
\begin{theorem}
    \[
        \mu_r' = \frac{d^r M_X(t)}{dt^r}|_{t=0} = M_X^{(r)}(0)
    \]
\end{theorem}
\begin{theorem}
    Let $a, b$ be constants. Then 
    \begin{enumerate}
        \item $M_{X+a}(t) = E(e^{(x+a)t}) = e^{at}M_X(t)$
        \item $M_{bX}(t) = E(e^{bXt}) = M_X(bt)$
        \item $M_{\frac{x+a}{b}}(t) = $
    \end{enumerate}
    
\end{theorem}


\section{Lecture 15 - 19/10/26}
\subsection{Product moments}
\begin{definition}
    The $r$th and $s$th product moment about the origin of random variables $X$ and $Y$ denoted by $\mu_{r,s}'$ is 
    \[
        \mu_{x,y}' = E(X^rY^s), \quad r, s = 0,1,2,3,\dots
    \]
    Recall that $X^rY^s$ is a random variable on $S$. and if $X, Y$ discrete,
    \[
        E(X^rY^s) = \sum_{x } \sum_{y} x^ry^sf(x,y) 
    \]
    for $f(x,y)$ joint pmf.
\end{definition}
\begin{definition}
    $r$th and $s$th product moment about the means of $X, Y$ is 
    \[
        u_{r,s} = E((X-\mu_X)^r(Y-\mu_Y)^s)
    \]
    where $\mu_i$ are the means respectively.
\end{definition}
\begin{definition}
    $\mu_{1,1}$ is called the covariance of $X$ and $Y$ denoted by $C[X,Y]$.
\end{definition}
\begin{theorem}
    \[
        C[X,Y] = E(XY) - E(X)E(Y) = \mu'_{1,1}-\mu_X\mu_Y
    \]
    
\end{theorem}
\begin{remark}
    The covariance measures the joint variability of $X$ and $Y$. If larger values of $X$ correspond to larger values of $Y$, then $C[X,Y] > 0$.
    If larger values of $X$ are more likely to correspond to smaller values of $Y$, then $C[X,Y] < 0$. 
\end{remark}
\subsection{Expectation of independent variables}
\begin{theorem}
    If $X,Y$ are independent, 
    \[
        E(XY) = E(X)E(Y), \quad C[X,Y] = 0
    \]
    
\end{theorem}

\section{Lecture 16 - 21/10/26}
\begin{theorem}
    If $X_1, X_2, \dots$ are pairwise independent, 
    \[
        Y = \sum_{i=1}^{n} a_i X_i 
    \]
    then 
    \[
        \sigma^2_Y = \sum_{i=1}^{n} a_i^2 \sigma^2_{X_i}
    \]
    And if $Y_1, Y_2$ defined above but for different constants,
    \[
        \sigma^2_{Y_1,Y_2} = \sum a_i b_i \sigma^2_{X_i}
    \]
    
\end{theorem}
\begin{theorem}
    Properties of covariance:
    \begin{enumerate}
        \item $\sigma^2_X = C[X,X]$
        \item $C[aX,Y] = aC[X,Y]$
        \item $C[X,Y] = C[Y,X]$
        \item $C[X+Y,Z] = C[X,Z] + C[Y,Z]$
    \end{enumerate}
    
\end{theorem}
\begin{theorem}
    \[
        \sigma^2_{X+Y} = \sigma_X^2+\sigma_Y^2+2C[X,Y]
    \]
     \[
        \sigma^2_{X-Y} = \sigma_X^2+\sigma_Y^2-2C[X,Y]
    \]   
\end{theorem}
\begin{definition}
    Correlation 
    \[
        P(X,Y) = \frac{C[X,Y]}{\sigma_X \sigma_Y}
    \]
    
\end{definition}
\subsection{Conditional expectation}
Suppose $X,Y$ are random variables and $Y=y$. What is the expetation/mean/variance of $X$?
The answer is they are called the conditional expectation/mean/variance of $X$ given $Y=y$.
\begin{definition}
    let $u(X)$ be a function of $X$. then 
    \[
        E(u(x),|y)
    \]
    is the conditional expectation of $u(x)$ given $Y=y$.
    For discrete,
    \[
        E(u(x)|y) = \sum_{x} u(x) f(x|y) 
    \]
    where $f(x|y)$ is the conditional pmf given $Y=y$. Same for continuous except with $f(x|y)$ the conditional pdf.
\end{definition}
\begin{definition}
    The conditional mean of $X$ given $Y=y$ is 
    \[
        \mu_{x|y} = E(X|y)
    \]
    
\end{definition}
\begin{definition}
    Conditional variance is 
    \[
        \sigma^2_{X|y} = E((x-\mu_{X|y})^2|y)
    \]
    
\end{definition}
\begin{remark}
    IMPORTANT.
    \[
        \sigma^2_{X|y} = E(X^2|y)-E(X|y)^2
    \]
    still holds for conditional variance.
\end{remark}

\subsection{Special discrete random variables}
Let $X$ be discrete and take values $x_1,x_2, \dots, x_k$. Assume
\[
    P(X=x_1) = P(X=x_2) = \cdots = P(x=x_k)
\]
Then the pmf is 
\[
    f(x_i) = P(X=x_i)
\]
and thus 
\[
    \sum f(x_i) = 1 \implies f(x_i) = \frac{1}{k}
\]
\begin{definition}
    A discrete random variable $X$ has a uniform distribution if 
    \[
        f(x) = \frac{1}{k}
    \]
    where $f(x)$ is the pmf of $X$
\end{definition}
\begin{definition}
    Let $X$ take values of either 0 or 1. Then $X$ has a Bernoulli distribution if its pmf is 
    \[
        f(x,\theta) = \theta^x (1-\theta)^{1-x},\quad x = 0,1
    \]
    
\end{definition}
\section{Lecture 17 - 26/10/26}
\begin{definition}
    A discrete random variable has a binomial distribution if its pmf is 
    \[
        b(x,n,\theta) = \binom{n}{x} \theta^x (1-\theta)^{n-x}
    \]
     
\end{definition}
\begin{theorem}
   \[
        b(x,n,\theta) = b(n-x,n,1-\theta)
   \]
\end{theorem}
Notice that it is hard ot get the mean and variance of $X$ with the definition of binomial distribution. Thus,
\begin{theorem}
    The moment generating function of the binomial distribution is
    \[
        M_X(t) = (1+\theta(e^t-1))^n
    \]
    
\end{theorem}
\begin{theorem}
    For the binomial distribution,
    \[
        \mu = n\theta, \quad \sigma^2  = n\theta(1-\theta)
    \]
    
\end{theorem}
\begin{remark}
    The Bernoulli distribution is a special case of the binomial distribution where $b(x,1,\theta)$.
\end{remark}
We want to further understand the binomial distribution. Our goal is to fix the mean and find an approxiamation for $b(x,n,\theta)$ where $n$ is very large, i.e. $n \rightarrow \infty$.

\begin{definition}
    For a discrete random variable $X$, it has a Poisson distribution if its pmf is 
    \[
        P(x,\lambda) = \frac{\lambda^x e^{-\lambda}}{x!}
    \]
    
\end{definition}

\begin{theorem}
    The moment generating function of the Poisson distribution is
    \[
        M_X(t) = e^{\lambda(e^t - 1)}     
    \]
    
\end{theorem}

\section{Lecture 18 - 28/10/26}
\subsection{Special pdf of continuous random variables}
\begin{definition}
    A continuous random variable $X$ has a uniform distribution if it has a pdf
    \[
        f(x,\alpha, \beta) = 
        \begin{cases}
            \frac{1}{\beta - \alpha}, \quad \alpha < x < \beta \\
             0, \quad \text{else}
        \end{cases}
    \]
    

\end{definition}
\begin{theorem}
    Assume $X$ has a uniform distribution. Then the mean is 
    \[
        \mu = \frac{\alpha + \beta}{2}
    \]
    and variance is
    \[
        \sigma^2 = \frac{1}{12} (\beta - \alpha)^2
    \]
        
\end{theorem}

\begin{definition}
    The gamma function is
    \[
        \Gamma(\alpha) = \int_{0}^{\infty} y^{\alpha-1}e^{-y}\, dy, \; \alpha > 0
    \]
    It holds the following properties:
    \begin{enumerate}
        \item $\alpha>1 \implies \Gamma(\alpha) = (\alpha-1)\Gamma(\alpha-1)$
        \item If $\alpha$ is a positive integer, then $\Gamma(\alpha) = (\alpha - 1)!$
        \item $\Gamma(\frac{1}{2}) = \sqrt{\pi}$
        \item For $n$ a positive integer, $\Gamma(n+\frac{1}{2}) = \frac{(2n-1)\cdot(2n-3)\cdots 3\cdot 1}{2^n} \sqrt{\pi}$
    \end{enumerate}
    
\end{definition}
\begin{definition}
    A continuous random variable $X$ has a Gamma distribution if it has a pdf
    \[
        g(x,\alpha, \beta) = 
        \begin{cases}
            \frac{1}{\beta^\alpha \Gamma(\alpha)} x^{\alpha-1}e^{\frac{-x}{\beta}}, \quad x > 0\\
            0, \quad x \le 0
             
        \end{cases}
    \]
    for positive $\alpha, \beta > 0$.
\end{definition}

\begin{definition}
    A continuous random variable $X$ has an exponential distribution if its pdf is 
    \[
        f(x, \theta) = g(x,1,\theta) = 
        \begin{cases}
            \frac{1}{\theta} e^{\frac{-x}{\theta}}, \quad x > 0 \\
             0, \quad x \le 0
        \end{cases}
    \]
    
\end{definition}
\begin{theorem}
    The $r$th moment about the origin of the Gamma distribution is 
    \[
        \mu_r' = \frac{\beta^r \Gamma(\alpha + r)}{\Gamma(\alpha)}
    \]
    
\end{theorem}
\begin{theorem}
    The mean and variance of the Gamma distribution are 
    \[
        \mu = \alpha \beta, \quad \sigma^2 = \alpha \beta^2
    \]
    
\end{theorem}
\begin{definition}
    $X$ has a normal distribution if its pdf is
    \[
        n(x,\mu, \sigma) = \frac{1}{\sqrt{2\pi}\sigma}e^{-\frac{1}{2}(\frac{x-\mu}{\sigma})^2}
    \]
    
\end{definition}

\begin{theorem}The following facts hold for the normal distribution:
    \begin{enumerate}
        \item $\int_{-\infty}^{\infty} n(x,\mu,\sigma)\, dx = 1$
        \item The moment generating function of a normal distribution is 
        \[
            M_X(t) = e^{\mu(t)+\frac{1}{2}\sigma^2t^2}
        \]
        
    \end{enumerate}
    
\end{theorem}

\begin{definition}
    Notation is $X \sim N(\mu, \sigma)$ means $X$ has a normal distribution with pdf $n(x,\mu,\sigma)$.
\end{definition}

\begin{theorem}
    For $X$ with a normal distribution, the mean is $\mu$ and the variance is $\sigma^2$.
\end{theorem}




\end{document}
